\chapter{Conclusion and Future Scope}\label{ch:conclusion}

\section{Conclusion}
We set out to build a system that treats the daily work of a Regional Transport Office as
a set of connected records and rules rather than as separate forms. The result is a
three-part application: a normalised MySQL database of 63 tables, a Java REST backend
that guards every rule inside a transaction, and a browser portal that gives staff and
citizens the same data in screens suited to their role.

Three ideas ran through the whole project. The first is that a rule worth having should be
enforced as close to the data as possible: unique keys, generated columns, CHECK rules and
triggers hold even if a program has a bug. The second is that concurrent use is the normal
case, so every operation that shares a scarce resource, such as an appointment slot, a
payment or a vehicle's current owner, takes its locks in a fixed order and writes
conditionally. The third is that a claim about quality should come with a measurement.
Accordingly the report gives the number of tests that pass (185 of 185), the load that was
sustained (about 380 requests per second with no errors at 20 users), and the effect of
an index (405.9~ms down to 1.4~ms on 500\,000 rows), together with the defects that the
testing exposed.

Along the way the work applied the content of five subjects of the course: object-oriented
programming in the layered Java code, database management in the schema and its
enhancements, software engineering in the requirements, design and testing process, data
structures in the maps, sets and trees used for transitions and regions, and the design and
analysis of algorithms in the ordering of locks, the cycle check and the complexity analysis
(Chapter~\ref{ch:subjects}).

\section{Objectives Revisited}
Chapter~\ref{ch:results} compared each objective with its evidence. In short, the schema,
the rule enforcement, the concurrency behaviour, the access control, the portal and the
measurements all reached the targets we set. The one objective with a qualification is
database-level enforcement: of the 21 CHECK rules defined, 19 are accepted by MySQL 8.0.44
and two (year of manufacture and payment time) are rejected because they use the current
date or time, which MySQL does not allow in such rules (defect D-08). That finding is
useful in its own right, and the correction is proposed in Chapter~\ref{ch:test}.

\section{Contributions}
\begin{itemize}
\item A complete relational design for an RTO, grouped into 13 modules, with ER diagrams, a data dictionary and a documented set of 25 indexes, 6 views, 5 triggers and 3 stored procedures.
\item A backend with 225 operations, permission checks loaded from the database, an audit trail written in the same transaction as the change, and notifications that are sent only after a commit.
\item A server-rendered portal of 70 templates with a menu that adapts to the permissions of the signed-in user, and friendly handling of errors and bad input.
\item An evidence set that other students can repeat: an integration-test suite, a schema audit, a crawl script, a load test and an index experiment.
\item A written record of ten real defects and their repair or proposed repair, which we hope is more useful to a reader than a list of passing results.
\end{itemize}

\section{Limitations}
\begin{itemize}
\item \textbf{External services are simulated.} Payments are recorded by gateway reference but no real gateway is connected; notifications are queued and marked sent without a real SMS or e-mail provider.
\item \textbf{Single-node design.} Revoked tokens are held in the memory of one backend process, and the portal keeps sessions in one process. Running several copies would need a shared store.
\item \textbf{Scale.} Load testing used one laptop and read-only requests; tables were tested up to 500\,000 rows, which is far below a national register.
\item \textbf{Two database scripts need correction.} The constraints script (D-08) and the SQL test script (D-09) are reported with proposed fixes but are not yet changed in the repository.
\item \textbf{Scope of the portal tests.} The portal was validated by a crawl of its pages and by manual walk-throughs rather than by automated browser tests.
\item \textbf{Paging.} Offset paging becomes slower for very deep pages on large tables.
\item \textbf{No documents are verified externally.} Uploaded files are checked for type and size only; they are not matched with any government record.
\end{itemize}

\section{Future Scope}
\begin{enumerate}
\item \textbf{Real integrations.} Connect a payment gateway with signed callbacks, an SMS and e-mail provider, and, where permitted, identity and vehicle registers.
\item \textbf{Shared infrastructure.} Move token revocation and sessions to a shared cache, so that the backend and the portal can run as several instances behind a load balancer.
\item \textbf{Automated browser tests.} Turn the crawl and the manual walk-throughs into repeatable end-to-end tests that run in continuous integration.
\item \textbf{Fix the database scripts.} Replace the two non-deterministic CHECK rules with a fixed bound and a trigger, update the SQL test script, and add the whole database build to the automated pipeline.
\item \textbf{Keyset paging and caching.} Use keyset paging for large lists and cache reference data that rarely changes.
\item \textbf{Mobile and public interfaces.} Because the portal and the backend are separate, a mobile application or a public status-tracking page can reuse the same API.
\item \textbf{Analytics.} Use the audit trail and the views for reports such as processing time per office, challan recovery rate and expiry forecasts.
\item \textbf{Accessibility and languages.} Add regional-language labels and an accessibility review of the portal.
\item \textbf{Document workflow.} Add scanned-document storage with virus scanning and a reviewer queue.
\end{enumerate}

\section{Lessons Learnt}
Building the system taught us that most of the effort in a data-heavy application sits in the
edges: what happens when two people click at once, when a date is in the future, when a
list page is requested with page number zero, or when a seed row points to a record that
is not there. Testing against a real database, and not a mock, found problems that no
amount of reading the code would have shown. We also learnt that a clear layering, with
validation, rules and storage each in their own place, makes defects easier to find, because
the layer at which a symptom appears usually points to the layer that must change. Finally,
writing the report forced us to check every figure against a run, and in several cases
that check corrected a number we had believed.
